\section{Hạn chế của khoa học}

\ 

Hiển nhiên rằng, không phải chỉ nguyên khoa học là bộ môn duy nhất đưa việc phản ánh thế giới làm trọng tâm. Thiên nhiên và con người liên tực được phản chiếu trong các bài văn thơ, bài hát hay những bức tranh. Hơn thế nữa, trong những tác phẩm nghệ thuật đó, ngoài việc chúng ta được cung cấp không gian và thời gian của sự vật, chúng ta còn nhìn thấy suy nghĩ nội tâm của các nhân vật, đặt mình vào trong dòng đời của người kể hoặc người được kể, để từ đó, chúng ta có thể trải nghiệm những khía cạnh của cảm xúc mà có thể chúng ta chưa từng được biết. Trái lại, tác phẩm khoa học thì khó có thể bộc bạch tinh thần của con người như vậy. Tác giả thường được nghe rằng cần phải trình bày các nghiên cứu như một ``câu chuyện khoa học'', nhưng tác giả cảm thấy câu chuyện này tương đối nhạt. Về mặt kết cấu, các tác phẩm theo văn phong khoa học vẫn có một kết cấu giống như các bài tự sự: có mở đầu, có bối cảnh, có hành động, có kết quả. Tuy nhiên, về mặt nội dung, việc thể hiện cảm xúc chủ quan là một điều cấm kị, do biểu cảm là thứ không cân đo đong đếm được. Nhiệm vụ của một người viết ra khám phá là phải làm sao cho người đọc, với một lượng kiến thức chuyên môn sâu nhưng không đến trình độ chuyên gia, có thể nhanh chóng hiểu và kiểm tra được tính chính xác của bài viết. Điều này dẫn đến sự thiếu hụt của tính ``người'', biểu qua sự hạn chế tới mức tối đa của các ngôn ngữ biểu cảm và các biện pháp tu từ trong các công trình khoa học được viết ra. Thực vậy, các sách sử học có thể chứng minh ảnh hưởng không mấy tích cực của phát xít Đức lên trẻ em sống ở các vùng bị đánh chiếm, nhưng liệu các số liệu và sự kiện trong những cuốn sách đó có tác động mạnh mẽ đến trái tim người đọc như các tác phẩm \emph{``Buổi học cuối cùng''} của An-phông-xơ Đô-đê\footnote{Alphonse Daudet (1840 -- 1897).} hay \emph{``Kẻ trộm sách''} của Mác-kớt Du-sắc\footnote{Markus Zusak (sinh năm 1975).} được hay không? Nếu bạn đọc chưa từng tiếp cận những tác phẩm này, thì hãy tìm đọc chúng, và tự đưa ra câu trả lời cho mình.